\documentclass[fontset=fandol, a4paper, 12pt, twoside]{ctexart}

% ==========================================
% 宏包加载
% ==========================================
\usepackage{geometry}       % 页面布局
\usepackage{fancyhdr}       % 页眉页脚
\usepackage{titlesec}       % 标题格式定制
\usepackage{setspace}       % 行距
\usepackage{xcolor}         % 颜色
\usepackage{marginnote}     % 边注处理

% ==========================================
% 页面布局设置 (宽边距批注模式)
% ==========================================
\geometry{
    a4paper,
    inner=2.5cm,      % 装订侧
    outer=6.5cm,      % 批注侧 (留出巨大空白)
    top=3cm, 
    bottom=3cm,
    marginparwidth=5cm, % 边注栏宽度
    marginparsep=0.8cm  % 正文与边注的距离
}

% 行距设置
\linespread{1.5}

% ==========================================
% 样式定制
% ==========================================

% 1. 页眉页脚
\pagestyle{fancy}
\fancyhf{}
% 偶数页：页码在左，文章标题在右
\fancyhead[LE]{\small \thepage}
\fancyhead[RE]{\small \textit{权力的网眼 / 米歇尔·福柯}}
% 奇数页：页码在右，当前章节在左
\fancyhead[RO]{\small \thepage}
\fancyhead[LO]{\small \textit{\leftmark}}
\renewcommand{\headrulewidth}{0.5pt}

% 2. 标题格式
\ctexset{
    section = {
        format = \Large\bfseries,
        number = \arabic{section},
        aftername = \hspace{0.5em},
        beforeskip = 20pt,
        afterskip = 15pt
    },
    subsection = {
        format = \large\bfseries\color{darkgray},
        number = {}, % 不带编号的小标题，保持流畅
        beforeskip = 15pt,
        afterskip = 10pt
    }
}

% 3. 引用环境 (福柯风格)
\renewenvironment{quote}
  {\list{}{\rightmargin=0cm \leftmargin=2em}%
   \item\relax\small\kaishu\color{black}}
  {\endlist}

% 4. 法语术语专用命令 (斜体)
\newcommand{\fr}[1]{\textit{#1}}

% 5. 边注样式 (用于提示核心概念)
\newcommand{\keyterm}[1]{%
    \marginpar{\raggedright\footnotesize\kaishu\color{blue} #1}%
}

% ==========================================
% 文档正文
% ==========================================
\begin{document}

% --- 标题页 ---
% \begin{titlepage}
%     \centering
%     \vspace*{4cm}
%     {\Huge \bfseries 权力的网眼}\\[1cm]
%     {\LARGE \textit{Les mailles du pouvoir}}\\[3cm]
    
%     {\Large \textbf{【法】米歇尔·福柯}}\\[0.5cm]
%     {\large (Michel Foucault)}\\[5cm]
    
%     \begin{quote}
%     \centering
%     选自《言与文》(Dits et Ecrits) 第315篇
%     \end{quote}
% \end{titlepage}

% \newpage

% % --- 正文 ---
% \setcounter{page}{1}

我们将着手尝试对权力概念展开分析。我远非是第一个试图规避弗洛伊德图式的人。这种图式使本能对立于压抑以及文化。在十多年前，整个精神分析学派都曾试图调整、构思这种本能对文化，本能对压抑的弗洛伊德图式。我所说的是诸如梅兰妮·克莱因（Melanie Klein）、唐纳德·威尼科特（Donald Winnicott）、雅克·拉康（Jacques Lacan）在内的英语圈以及法语圈的精神分析师们。他们已经试图表明，压抑远非是一种试图控制天赋的本能游戏的次级、外部（\fr{ultérieur}）且晚生（\fr{tardif}）的机制。压抑是本能机制的一部分，或者说，至少是性本能作为驱力发展、展开、形成自身的过程的一部分。

弗洛伊德的Trieb（驱力）概念不应该被解释为一种简单的自然给定（\fr{donnée naturelle}），一种被压抑施加了禁令之法的自然的生物学机制。根据精神分析师们的意见，Trieb应该被视作压抑深深渗入其中的某物。欲望（\fr{désir}）在需求（\fr{besoin}）、阉割、匮乏、禁令、法律（\fr{loi}）这些要素中被构成性欲（\fr{désir sexuel}）的。而这意味着弗洛伊德在十九世纪末所构想的性本能的原始概念发生的转变。因此，人们不应该将本能思索为一种自然给定，而应该总是某种阐释的整体，总是某种身体与法律、身体与确保对人群进行控制的文化机制之间的复杂游戏。
因此，我认为精神分析师们已经通过提出一种新的本能概念（\fr{notion}），或者说崭新的本能、驱力、欲望的概念架构（\fr{conception}，译注：此处仅为区别于\fr{notion}，后文中仍作概念），从而深刻地变更了这一问题。然而，令我困惑，或者至少让我不够满意的是，虽然精神分析师们在他们提出的这种阐释中可能修改了欲望的概念（\fr{conception}），但是他们绝对没有修改权力概念。

精神分析师们仍旧将权力的所指、权力的核心乃至权力所涉及的视作是禁令，是法律，是说“不”的事实（\fr{le fait de dire non}），是“你不应”这一公式和形式。在本质上，权力就是去说“你不应”。在我看来，这种权力概念——我将之后再谈——整体而言是不充分的，它是一种法律概念，一种权力的形式概念。我们应该促成另外一种权力概念，一种能够使我们更好地理解西方社会中权力与性态关系的概念。

我将试着发展，更准确地说，展示出我们能够在什么样的方向上发展出一种权力分析。这种权力并非单纯只是一种法律的，否定性的概念，而是作为一种权力技术的概念。

我们经常在精神分析师、心理学家与社会学家们看到这样一种权力概念。根据这种概念，权力本质上乃是规则、法律、禁令,是去标记被许可之物与被禁止之物之间的界限。我认为这种权力概念在十九世纪末，已经经由民族学（\fr{ethnologie}）精辟地提出并广泛地发展起来了。民族学总是热衷在与我们所处社会不同的社会中，探查作为规则体系的各种权力体系。而当我们自己试图对我们的社会及其权力行使方式进行反思时，我们在本质上是基于一种法律概念进行反思的——权力在哪里？谁占有了权力？哪些规则支配着权力？权力将哪种法律体系建立在社会体之上？

因此，围绕着我们的社会，我们总是在进行一种权力的法社会学（\fr{sociologie juridique}）研究。并且，当我们研究与我们社会相异的其他社会时，我们所进行的民族学研究在本质上也是一种规则与禁令的民族学。例如，我们可以看到，在从涂尔干到列维·斯特劳斯的民族学研究中，反复出现并且持续被阐释的问题乃是禁令的问题，或者说在本质上是乱伦禁令的问题。进而，从这一母题、这一本质上是乱伦禁令的问题内核出发，人们试图把握体系的一般运作。一直到最近几年，我们才看到一些关于权力的新观点的出现。这些观点要么是严格的马克思主义的，要么是远离经典马克思主义的。总之，以皮埃尔·克拉斯特（Pierre Clastres）\footnote{参考皮埃尔·克拉斯特收录于《社会反对国家：政治人类学研究》中的研究。\fr{La Société contre l'État. Rechercher d'anthropologie politique}, Paris, Éd. de Minuit, coll. «Critique», 1974.}的研究为例，我们看到一种将权力视作技术的，崭新的权力概念出现了。这种新的权力概念试图将将自己从规则与禁令的首席权或者特权中解脱出来——后者实际上自涂尔干起一直到列维·斯特劳斯，都一直支配着民族学研究。

总之，我想提出的问题乃是我们的社会，或者一般而言的西方社会，是如何会以一种如此限制，贫困并且否定的方式来构思权力的？为什么我们总是将权力作为法律，或者禁令来加以构思？为什么法律和禁令这一模式会有如此特权？很明显，我们可以认为这一情况源自于康德的影响，源自于一种理念。在这种理念的最新范例中，道德法则，或者说“你不应”的律令，乃至于“你应该”/“你不应”这一对立，都是对人类行为所进行的一切调节的母题。但老实说，这种诉诸康德的影响来进行的解释，很明显在总体上是不充分的。问题的关在在于知道康德是否有过这种影响，以及为什么能够有如此强的影响。为什么涂尔干这位在法国第三共和国初期带有模糊的社会主义色彩的哲学家，在决定分析某一社会的权力机制时，会则以这样一种方式倚仗于康德？

我认为，我们可以通过以下的术语来对上述的状况进行一番简略分析。实际上在西方，自中世纪以来，各种被建立起来的重要体系，都是借助于君主权力（\fr{pouvoir monarchique}）增长发展起来的。而其代价乃是权力，或者更精确地说，各种封建权力（\fr{pouvoir féodal}）。而在这种君主权力与封建权力的斗争中，法（\fr{droit}）既是君主权力对抗各种其他团体（\fr{institution}）、习惯（\fr{moeur}）、规约（\fr{règlement}）的工具，也是君主权力与各种封建社会的特有的隶属和束缚形式抗衡的工具。我将简要展示两个例子。一方面，西方世界的君主权力的发展在很大程度上依赖于各种司法机构的发展，同时，其自身也通过发展这些机构得以扩张。通过内战，君主权力终于用一整套的法律和法院体系，取代了古老的私人纷争解决办法。而这套系统在事实上赋予了君主权力以解决各种个人纠纷的可能性。同样，十三、十四世纪在西方再次出现的罗马法，在君主们的手中，以牺牲封建权力为代价，成为了确定君主们自身权力之形式与机制的有效工具。另一方面在欧洲，国家的发展部分地被确定为，或者说被作为一种法律思想的发展。君主权力，国家权力本质上都被表象于法之中。

与此同时，市民阶层（\fr{bourgeoisie}，译注：后文中在更近代以及资本主义的语境里我们会将其转而翻译为布尔乔亚）也利用了王权（\fr{pouvoir royal}）的扩张以及封建体系的衰退和萎缩。市民阶层乐于推动法权利（\fr{droit}）体系的发展。这种法权体系允许，或者说能够给予市民阶层保证其社会发展之所需的经济交换框架。因此，法权利的词汇和形式成为了市民阶层与君主所共有的权力的表象/代表体系。自中世纪末期一直到十八世纪，市民阶层与君主逐渐成功建立起了一种权力形式。这种权力形式被表现，或者被给定为法权利词汇的话语和语言。当市民阶层最终开始摆脱君主权时，他们使用的正是这种话语——尽管这种法律话语曾经属于君主制，但如今却被市民阶层转而用于对抗君主制自身。

我们来看一个简单的例子。当卢梭构筑其国家理论时，他试图展现的是主权者如何诞生的——虽然这是一位集体主权者（\fr{souverain collectif}），是作为社会体的主权者，或者更准确地说，是作为主权者的社会体。他想展现这位主权者诞生于个体权利的出让（\fr{cession}）与转让（\fr{aliénation}），诞生于每一个体都有义务去承认的那些禁令法则的公布。他们之所以有这一义务，乃是因为他们自己就是主权者的一员，他们自己就是主权者，因此，是他们自己给自己施加了这些法律。总之，人们通过这一理论机制对君主制制度进行批判。而这一理论工具正是法权利这一曾经由君主制自己所创造的工具。换言之，除了法权利，除了法律体系，西方从未有过另外一种表象/代表体系，另外一种权力分析或者权力公式。我认为，正是这个原因最终导致我们直到最近都无法在不使用法律、规则、主权、授权等等本质的、根本的（诸如此类）的概念的前提下进行权力分析。我认为，如果我们想要推进一种不再是对权力的表象/代表，而是对权力真正的运作展开的分析，那么，我们如今应该摆脱的，正是这种基于法律与主权者，规则与禁令的，法律意义上的权力概念。

那么，我们如何能够在权力的肯定性（\fr{positif}）机制中，尝试纪念性权力分析呢？我认为，我们可以在某些文本中发现一些进行这类分析的基本要素。我们或许可以在边沁这位十八世纪末的和十九世纪初的英国哲学家这里发现这些要素。实际上，边沁乃是市民/资产阶级权力（\fr{pouvoir bourgeois}）的伟大理论家。同样，我们也可以在马克思那里，特别是在《资本论》第二卷中发现这些要素。我认为正是在这些地方，我们能看发现一些有助于我们在权力的肯定性机制中进行权力分析的要素。

总体而言，我们首先在《资本论》第二卷中所能看到的是，并不存在一个权力，而是存在许多权力\footnote{Marx (K.), Das Kapital. Kritik der politischen Ökonomie. Buch II : «Der Zirkulationsprozess des Kapitals», Hambourg, O. Meissner, 1867. (\fr{Le Capital. Critique de l'économie politique}, livre II : «\fr{Le procès de circulation du capital}», trad. E. Cogniot, C. Cohen-Solal et G. Badia, Paris, Éditions sociales, 1976, vol. II).}。许多权力指的是，存在着多种支配形式，多种臣服形式。这些形式局部地发挥作用——例如，在工厂车间中，在军队里，在奴隶制的产业，或者在在有奴役性关系的产业中。所有这一切，都是权力的局部、区域性的形式。它们有着各自的运作模式、程序和技术。所有这些权力形式都是异质的。如果我们先进行一种权力分析，我们不能讨论唯一的权力（\fr{le pouvoir}），我们应该谈论各种权力（\fr{des pouvoirs}），情切应该尝试在其特殊的历史与地理之中对它们进行局部化。

一个社会不是一个只有一种且唯一种权力运作于其中的一元体（\fr{corps unitaire}）。实际上，社会是诸多仍旧保持其特殊性的权力所构成的并置（\fr{juxtaposition}），关联（\fr{liaison}），协调（\fr{coordination}），乃至等级秩序。例如，马克思非常强调工厂主在工厂车间中行使在权力相对于在社会其他领域中运作的法律型权力所具有的特殊性、相对自主性，以及某种程度上的不可穿透性。因此，存在着不同的权力区域（\fr{region de pouvoir}）。社会是由各种不同权力构成的群岛。

其次，权力似乎不能且不应该被简单理解为占据优先地位的中央权力的分支或者结果。诸如格劳修斯（Grotius），普芬多夫（Pufendorf）或卢梭等思想家的法学家图式所要谈及的正是“起初，社会并不存在。随后，当某一个主权中心点（其组织起了社会体）出现时，社会出现了。而这个主权中心点也使得一系列局部和区域性的权力成为可能。”而马克思暗中却不接受这一图式。相反，马克思指出大型的国家机器是如何产生于各种各样先在且原初的较小的权力区域的——这些权力区域可以是财产权，可以是奴隶制，可以是车间厂房也可以是军队。实际上，国家的统一性相较于先在的那些特殊且区域性的权力而言是次级的。

第三，这些特殊的、区域性的权力并不绝对地具有一种进行禁止、禁锢，一种表达“你不应”的优先功能。实际上，这些局部且区域的权力的原始、本质且持续的功能乃是去成为某种效力（\fr{efficience}）或资质（\fr{aptitude}）的生产者，成为某种产品的生产者。例如，马克思就曾经对工厂车间与军队种的纪律问题进行过非常精彩的分析。虽然接下来我想针对军队中的纪律的分析不见于马克思的文本中，但这不重要。自十六世纪末、十七世纪初起，直到十八世纪末为止，军队之中发生了什么？当时的军队在根本上是由许多围绕在军事首脑身边组织起来的，相对可以互换的个体统一体所组成的。此后，这种组织模式开始被一种巨大的金字塔结构的统一秩序所取代。这一金字塔结构包含着和一系列士官或者技术人员在内的中间首脑。而这一切本质上是源于当时军事人员所进行的一项技术创新，即能够相对更快而且更密集地开火的步枪。

自此之后，人们不再能够将军队（让其运作起来是非常危险的）视作一些孤立的小统一体以及一些可互换的要素的组合。为了让军队更卓有成效、为了让人们能够更好地使用枪支、为了能保证够确保士兵在一条漫长的战线上站在其应该站的位置上、为了让士兵们能同时站成笔直的一条线列、等等等等，每个个体都应该被更好地操练起来。纪律的问题指明了一种与新的，士官们的权力技术，指明了由士官、中级军官和高级军官们共同组成的等级秩序。而通过根据每个人和每个位置的特殊性所共同构成的整体的统一，这种等级秩序使军队能够发挥出其最大效能。由此，军队便可以被视作是一种更复杂的等级制统一体。

得益于这种新的权力程序，这一时期的军事表现非常出色。而这种权力程序的功能也绝非是去禁止什么。当然，这种权力程序确实被用于禁止这或禁止那，但是，这种权力程序的目标绝非是去说“你不应”。其更本质的目标是去为军队塑造更好的表现、产品以及生产性。借助这种新的权力技术，军队不断地作为死亡的制造者趋于完善，或者说得到确保。而这种权力技术绝非是禁令。我们同样可以如此评价工厂车间中的纪律问题。这些纪律发源于十七世纪和十八世纪。在这个时期，拥有众多工人——几百人——的大型工厂开始代替那些行会性质的小工场。因此，人们就必须同时监视以及协调人们与他人、与其他劳动分工之间的动作姿态。同时，劳动分工也正是迫使人们发明出新的工厂纪律的原因所在。但是，我们反过来也能说，工厂纪律是人们能够进行劳动分工的条件。如果没有这种工厂纪律，换言之，如果没有等级秩序，没有监视，没有工头的出现，没有对动作姿态精确（\fr{chronométrique}）控制，人们是不可能进行劳动分工的。

最后是第四点重要看法。我们应该将这些权力机制、权力程序视作是一些技术，或者说一些程序。这些技术或者程序被发明出来、被加以完善，并不断得到发展。存在一种切实的权力技术，或者更准确地说，存在着各种切实的权力技术，这些权力技术有其自身的历史。关于这点，我们能够再次轻松地在《资本论》第二卷中发现一种分析，或者至少说是分析的草图。这种分析（或者分析的草图）勾勒了权力技术的历史，勾勒了在车间与工厂中被行使的权力技术的历史。我将在思考性态时，追随这些根本性的指示，并常识不从法律视角，而是从技术视角来定义权力。

实际上，在我看来，如果我们通过一种特权化国家机器的方式对权力进行分析，如果我们只是将权力视作一种保守（\fr{conservation}）机制，如果我们只是将权力视作是一种法律的上层建筑，那么我们所做的，只不过是重新提出布尔乔亚思想中的经典命题——这一命题在根本上将权力认定为一种法律事实。对国家机器、保守功能以及法律上层建筑这些面相的特权化，实际上是“卢梭化”（\fr{rousseauiser}）了马克思，是复述布尔乔亚式的权力理论或者法律式的权力理论而已。不足为奇的是，那种所谓的马克思主义的权力概念，即将权力视作是国家机器、保守功能的审判体系（\fr{instance}）、法律的上层建筑的那种所谓的马克思主义权力理论，基本上可以在十九世纪末的欧洲社民主义思潮中发现——当时的问题核心正是在于寻求能够使马克思在布尔乔亚法律体系内运作的方法。但是，我更希望藉由重提《资本论》第二卷中的这种权力分析来远离国家机器、权力再生产功能，法律上层建筑特征等一系列后来添加或改写的，具有特权性质的话题，从而试着去考察是否可能书写一部关于西方世界中各种权力的——或者更根本地说，各种被投入于性态领域中的权力的历史\footnote{1981年的公开出版版本止于此处。}。

那么，基于这一方法论原则，我们如何能够书写一部与性态相关的权力机制史呢？我认为——尽管某种程度上非常图式化——可以这么说。自中世纪晚期以来由君主制所成功组织起来的权力体系对于资本主义的发展有两大不利之处。第一，施加于社会体之上的政治权力是一种非常不连续的权力。这种权力之网的网眼过于巨大，数不尽的东西、要素、行为乃至程序都逃离了权力的控制。我们可以用例子来说得更清楚。我们注意到，直到十八世纪末，走私在整个欧洲都占据重要地位。这种经济流通和其他经济流通同等重要，同时还完全逃脱了权力管控。在另一个角度上，这也是人们的生存条件之一——如果没有海盗，商业将无法运作，而人们也就无法生存。换言之，非法活动（\fr{illégalisme}）曾经乃是生命条件之一，但其同时也意味着某周脱离于权力之外，权力无法控制的东西。总之，有一些经济过程，有一些机制在某种程度上外在于权力的控制。而这种情况就要求建立起连续、精确，某种意义上原子化的权力。从一种有间隙的、整体的权力，过渡到一种连续的、原子化和个体化的权力。在后者中，每一个东西，每一个显现在其身体之中，体现在其姿态之上的个体，本身都能够被控制。这样一种连续的权力，代替了整体和粗放的控制。

这种运行于君主制内部的权力机制的第二点重大不利在于，它的运作成本极为昂贵。而其昂贵的原因正在于这种权力的功能——这种权力所包含的功能——在本质上是一种汲取（\fr{prélèvement}）的权力，是一种具有拥有征收某物的权利或者力量的权力。征收的可能是捐税，也可能是教士们对成熟庄稼所征收的什一税——总之，是主人、王权或者教士们所征收的强制的或者百分比的税收。因此，这种权力本质上是征税官或掠夺者。在这种意义上，这种权力总是在经济上做减法，并且最终无法鼓励和刺激经济流通。因而，这种权力总是经济流通的障碍与约束。基于这第二项担忧，产生了第二种必要性，即发现一种新的权力机制的必要性。这种权力机制运行于经济过程的内部，可以在最大程度上控制各种人和物，但同时对社会而言即不昂贵，且本质上亦非掠夺。


通过这两个目标，我认为我们可以粗略地理解西方世界中的权力技术的重大转变。本着一种较为初级的马克思主义精神，我们习惯于认为，伟大的发明指的是蒸汽机或者诸如此类的东西。这是对的，并且这类发明也确实非常重要，但是还有另外一系列技术发明在重要程度上不逊色于这类发明，并且在根本上乃是其他发明得以发挥功效的先决条件。政治技术也同样如此。在整个十七和十八世纪，在权力形式上也有一整套新的发明。总之，我们不仅应该书写一部工业技术的历史，同样也应该书写一部政治技术的历史。并且，我认为我们能够将这些被我们特别标定在十七、十八世纪的权力技术的发明分为两大项。我之所以将它们分为两大项，是因为在我看来，这些技术是沿着两条不同的方向发展的。一方面，存在一种我称之为“纪律”的技术。实际上，纪律指的是一套权力机制，通过这套权力机制，我们可以控制社会体中哪怕最微小的要素，可以触及社会原子（换言之就是个体们）自身。纪律是权力的个体化技术。在我的用法上，纪律是指如何监视某人，控制其行为（\fr{conduite}）、举止（\fr{comportement}）、资质，强化其表现，拓展其能力，使其更有用的技术。

我刚才为各位举了军队纪律这一例子。这是一个非常重要的例子，因为军队确实是纪律这一伟大发明最初形成和发展的场所。而这项技术的发明与另一项技术-工业性质的发明，即能够更快进行射击的步枪相联系。实际上，自此以后我们可以说，士兵们不再是可互换的，不再是单纯简单的炮灰，也不再简单只是个能打之人。为了成为一名称职的士兵，人们应该懂得如何射击，为此人们应当首先度过一个练习阶段。士兵同样需要知道该如何移动，如何协调自己与其他士兵的姿势，总之，士兵们变成了某种具备技巧的东西（\fr{chose}），因此也就是某种珍贵的东西。一旦士兵变得珍贵，那么就越要保护他们，而为了保护他们，就越有必要教会他们在战斗中保命的技巧，而越是要教授给他们这些技巧，那么练习阶段就越长，士兵也就越珍贵。突然之间，军事训练技术有了长足发展，并在腓特烈二世麾下著名的普鲁士军队中达到了顶点——腓特烈的军队花费了他们生命中主要时间在操练之上。准确来说，普鲁士军队这一普鲁士纪律的楷模，乃是士兵身体纪律的最大化完善和增强，并且在某种程度上乃是其他纪律的典范。

而另一个我们借此观察到新的纪律技术出现的领域则是教育领域。这种纪律方法首先出现在中学里，随后拓展到了小学中。在这些地方，个体们都在其杂多性（\fr{multiplicité}）被加以个体化。学校中聚集了几十，成百，有时候甚至上千的中学生、小学生。而施加在他们之上的权力，其成本应当远远低于家庭教师所具有的那种只能存在于一位老师与一位学生之间的权力。在学校中，我们可以让一位老师配对十几名学生。然而，尽管学生杂多，人们却得到了一直权力的个体化，一种持续的控制以及一种对每一个瞬间的监视。在这个空间中，这样一种所有上过中学的人都很熟悉人物出现了。其作为监视者，位置在金字塔结构中与部队里的士官是相符的。而量化分级的概念、考试、会考也都出现了，最终，在老师的眼中，或者通过施以每个学生的资格或判断，人们可能可以给每个学生划定出其应该在的位置。


各位试想一下，现在大家排成排坐在我对面前。这种坐法对大家来说可能非常自然，但是我们有必要记住，这种坐法的出现在文明史上其实是相对晚近的事情。在十九世纪初，我们仍然可以在一些学校看到学生们三五成群站在授课老师身边的场景。很明显，这表明了教师无法切实且个体地监视他们——这里有一组的学生，然后才有一位教师。而如今，学生们被有序排好，教师的视线可以个体化每一个人，可以通过点名来知道他们是否到场，在干什么，是否清醒，是否在打哈欠等等。这都是些琐事，但是恰恰是些非常重要的琐事。因为最终，在权力的一系列的行使之中，正是这些小小的技术使得新的权力机制得以投入应用，得以发挥功能。在整个十九世纪中，我们也可以在工厂车间内看到这些发生在军队或者学校里的事情。而这就是我所说的权力的个体化技术——一种紧盯个体身体、行为的技术，大致上，我所说的是一种政治解剖学（\fr{politique anatomie}），一种解剖-政治（\fr{antomo-politique}），一种解剖个人的解剖学。

这就是出现在十七和十八世纪的一组权力技术，而在稍晚的十八世纪后半叶，我们则可以看到另一组权力技术的诞生。这组权力技术尤其是在英国得到了发展（有必要指出，令法国羞愧的是，前一种技术主要发展于法国和德国）。这组技术的目标不是那些作为个体的个体，相反，是人口。换言之，人口这一重要的东西在十八世纪已经被发现了。在这种权力技术中，权力不是简单地被行使于各个臣民（\fr{sujet}）之上（这是君主制的核心主题，根据这一主题，总是有一位主权者和许多臣民）。我们可以看到，权力被行使于人口之上。那么，人口指的是什么呢？人口并不是简单指说一个大型的人类团体，而是一些被某些生物学法则或者过程渗透、指导或者规定的活存在（\fr{être vivant}）。一组人口有其出生率、死亡率、年龄曲线、年龄金字塔、发病率、健康状态。一组人口可能会走向消亡，或者相反地走向繁荣。

这一切都是在十八世纪开始被发现的。总之，我们注意到权力与主体的关系，或者更准确的，权力与个体的关系的形式，不仅仅只是臣服形式（\fr{forme de sujétion}）这一种。臣服形式允许权力从臣民那里征收财货、财富乃至他们的身体和血。如果我们想要更精确地使用“人口”概念，即将其看作是一种为了生产诸如财货、财富，乃至生产其他个体的机器，那么在我们所讨论的这种关系中，权力应当行使于那些组成一种必须被加以考虑的生物学实体的个体之上。发现人口同时也是发现个体以及发现可训练的身体。后者乃是西方政治进程在其变化过程中的另一组重要的技术核心。正是在这一时期，相对于我前面提到解剖-政治，人们发明了我称之为生命-政治（\fr{bio-politique}）也正是在这一时期，我们看到了诸如居住条件、城市生活条件、公共卫生、出生率与死亡率关系变化等问题的出现。也正是在这一时期，产生了诸如我们应该如何让人们多生小孩、如何规制人口流动、如何规制人口增长率、移民增长率等问题。也正是从那时候起，所有这一系列的观察技术——很明显包括同统计学，除此之外是大规模的行政、经济、政治机关——都承担着调控人口的职责。在权力技术中有两场大革命，即纪律的发明与调控的发明，也就是一种解剖-政治的完善化与一种生命-政治的完善化。

自十八世纪以来，生命就成为了权力的对象。生命和身体都成为了权力的对象。此前，只存在着一些主体/臣民，一些人们可以从其身上汲取财货、生命乃至其他东西的法律主体。如今，则存在着许多的身体与人口。权力成为了唯物论者。权力在本质上不再是法律的了。权力必须要处理诸如身体、生命这样一些实在的东西（\fr{chose réelle}）了。生命开始跻身于权力的领域当中——而这是一个重大变化，或许是人类社会历史上最重要的变化之一。我们能够很清楚地看到，从这一时期起（换言之，正是自十八世纪以后），性（\fr{sexe}）是如何能够成为一个绝对重要的一个环节。实际上，这是因为性被非常精准地置于将对身体的个体性纪律与对人口的调控啮合起来的点位上。通过性，我们能够确保对个体的监视。同样，我们也知道，这就是为什么在十八世纪，而且正是在中学里，青少年们的性态会成为一项医学问题、道德问题，甚至近乎成为具有头等重要性的政治问题。这是因为，通过对性态的控制——同时也是以控制性态为托辞，人们可以监视那些中学生、青少年的生活，监视他们的每一个瞬间——甚至在他们睡觉时也能监视他们。性由此成为了“纪律化”（\fr{disciplinarisation}）的工具。很快，性将成为我所说的解剖-政治中最为本质性的要素之一。但是，在另一方面，也正是性，确保了人口的再生产，正是通过性，正是通过一种关于性的政治，人们能够转变出生率与死亡率之间的关系。无论如何，性的政治都将被整合进整个关于生命的政治（其在十九世纪会变得格外重要）内部。性正位于解剖-政治与生命-政治的啮合处之上，其乃是纪律与调控的交汇处。在十九世纪末，为了使社会成为一架生产机器，性在其已经具有的功能之中，成为了具有首要重要性的政治环节。

\end{document}